\subsubsection{Mathetmatical explanation for the dip in the magnitude of force field}
Consider a simple potential formulation with a skewed distance.
\[
U = \frac{1}{k}, \qquad k = \sqrt{x^2 + \frac{y^2}{\gamma^2}}
\]

\[
\frac{\partial U}{\partial x} = -xk^{-3}, \qquad
\frac{\partial U}{\partial y} = -\frac{y}{\gamma^2}k^{-3}
\]
\[
|\nabla U|^2 = \frac{x^2 + y^2/\gamma^4}{k^6}
= \frac{x^2 + y^2/\gamma^4}{(x^2+y^2/\gamma^2)^3}
\]

We fix $y=c$ and if the force field has no dip then $|\nabla U|^2_{y=c}$ should have a local maximum at $x=0$
\[
g(x) = |\nabla U|^2\Big|_{y=c} = \frac{x^2 + c^2/\gamma^4}{(x^2+c^2/\gamma^2)^3}
\]
\[
g'(x) = \frac{2x}{(x^2+c^2/\gamma^2)^3} + \frac{(x^2+c^2/\gamma^4)(2x)(-3)}{(x^2+c^2/\gamma^2)^4}
\]

Setting $g'(x)=0$:
\[
x^2 = \frac{3(x^2+c^2/\gamma^4)}{(x^2+c^2/\gamma^2)}
\implies 2x^2 + \frac{c^2}{\gamma^2} = 3x^2 + \frac{3c^2}{\gamma^4}
\]
\[
x^2 = c^2\left(\frac{1}{\gamma^2} - \frac{3}{\gamma^4}\right)
\]

For no dip, no real solution should exist:
\[
\frac{1}{\gamma^2} \le \frac{3}{\gamma^4} \implies \gamma \le \sqrt{3}
\]
So whenever skewness ratio $\gamma\ge \sqrt{3}$ the force field has a dip at $x=0$.

\subsubsection{Intuitive explanation for the dip}


Alternatively, using $|\nabla k|$:
\[
\frac{\partial k}{\partial x} = \frac{x}{k}, \qquad
\frac{\partial k}{\partial y} = \frac{y}{k\gamma^2}
\]
\[
|\nabla k| = \frac{1}{k}\sqrt{x^2 + \frac{y^2}{\gamma^4}}
\]
\[
F = \frac{|\nabla U|^2}{|\nabla k|^2} = \frac{1}{k(x^2+y^2/\gamma^2)^3}
\]

Fixing $y=c$ ($c \ge 0$):
\[
F = \frac{1}{k(x^2+c)^3}
\]

$F$ is maximum at $x=0$.